\section{Service-Matched Delivery: What HACT Can Save}\label{sec:delivery}
\subsection{Three consumers define three services}
The local aggregator receives bounded packets; a logging collector may receive their complete stream; a supervisor normally receives a status projection and retrieves details selectively. Define $\psi(e,c)$ as canonical status text containing semantic identity, counts and local authorization, but no display positions. For the same trusted semantic report,
\begin{equation}\label{eq:rootcontext}
 \psi_{L_1}(e,c)=\psi_{L_2}(e,c).
\end{equation}
Its byte length and token count under any fixed deterministic tokenizer are therefore identical. This is a useful negative control: layout optimization cannot shrink a layout-independent root input. A flat ledger similarly provides final per-check states without preserving the complete intermediate hierarchy. HACT's layout objective applies when the hierarchical diagnostic stream itself is part of the required service.

Diagnostics preserve full header/card records and page whole packets under a 16,384-byte limit. A packet that cannot fit is rejected rather than truncated. For a real supplied tokenizer $\tau$ and explicitly reserved framing budget $r$, token-aware admission additionally requires
\begin{equation}
 |p|_{\mathrm{UTF8}}\le C_{\mathrm{bytes}},\qquad
 |\tau(\mathrm{decode}_{\mathrm{UTF8}}(p))|+r\le C_{\mathrm{tokens}}.
\end{equation}
The implementation retokenizes each complete proposed page, including its prefix; it does not assume additivity of fragment token counts. The optional measurement command requires an actual \texttt{tiktoken} encoding~\cite{tiktoken}. Provider framing and billed usage remain properties of a provider adapter, not of the byte objective.

\subsection{Codec boundaries are part of the service}
Canonical unpadded JSON, independently compressed packets, and whole-epoch DEFLATE preserve the same logical records but expose different streaming and compression behavior. The padded objective is therefore not claimed to optimize compressed bytes or model tokens exactly. Practical selection is advisory until the deployed representation is measured with the deployed codec.

This distinction also separates representation choice from verification choice. If a deployment already requires protected hierarchical verification, fixed and learned layouts can be compared at marginal cost under the same service. If the baseline used ordinary testing with no protected evidence service, adopting the verifier is a separate engineering decision and its additional cost cannot be charged to, or justified by, layout compression alone.

\subsection{Break-even accounting}\label{sec:breakeven}
Consider $q$ copies of the \emph{same} workload and delivery service through one nonoverlapped bottleneck. Let $B_0,B_1$ be codec-matched application bytes per copy, $R>0$ effective application goodput, $E\ge0$ additional encoding/decoding work per copy, and $C\ge0$ consistently amortized one-time incremental work. Common setup cancels only when it is actually common. The conditional net benefit is
\begin{equation}\label{eq:latency}
 d=(B_0-B_1)/R-E,\qquad \Delta(q)=qd-C.
\end{equation}
A strictly profitable integer copy count exists only when $d>0$, and then
\begin{equation}\label{eq:qmin}
 q_{\min}=\lfloor C/d\rfloor+1.
\end{equation}
A zero gain retains the baseline. The implementation uses decimal arithmetic at equality, rejects mismatched service IDs, and returns \texttt{UNKNOWN} when required measurements are absent. An explicit zero is accepted only as a declared assumption or observation.

When both variants use the same protected gate, $C$ represents extra layout work rather than the entire gate cost. Repeating delivery can amortize preparation of the same candidate; it cannot amortize preparation of future candidates as though they were the same workload. Parallel links, computation overlap, cached diagnostics, queueing, and connection setup require direct critical-path measurement. RTT alone does not increase the byte-dependent term when both alternatives pay the same setup and request count.

\subsection{Remote-measurement boundary}
The artifact includes an opt-in TLS measurement kit with certificate validation, hostname checking, bounded frames, validated decompression, and acknowledgments that bind received length and digest. Local tests exercise cold and reused sessions, rejected hostnames, untrusted roots, malformed acknowledgments, and truncated records. These tests validate the measurement path; they are not WAN performance measurements.

A qualifying remote experiment must record host/path provenance, codec, connection policy, matched payloads, application goodput, and transport observations such as RTT or retransmissions where available, following the separation between application payload and transport behavior emphasized in TCP throughput testing~\cite{rfc6349}. TLS authenticates the configured server, not the truthfulness of a checker. Attestation, Byzantine consensus, client authorization, non-repudiation, and OS isolation remain outside HACT's local semantic guarantee.
